% Intended for Inverse Problems (IOP Publishing).
% main.tex prefers the current official iopjournal.cls when it is present.
% The bundled iopjournal-draft.cls is only a local compilation fallback.
\IfFileExists{iopjournal.cls}{\documentclass{iopjournal}}{\documentclass{iopjournal-draft}}

\usepackage[T1]{fontenc}
\usepackage{graphicx}
\usepackage{xcolor}
\usepackage{hyperref}
\usepackage[utf8]{inputenc}
\usepackage{lmodern}
\usepackage{amsmath,amssymb,amsthm,mathtools}
\usepackage{bm}
\usepackage{microtype}
\usepackage[numbers,sort&compress]{natbib}
\usepackage{xurl}
\usepackage{booktabs}
\usepackage{enumitem}
\usepackage{listings}
\usepackage{placeins}

\newtheorem{theorem}{Theorem}[section]
\newtheorem{proposition}[theorem]{Proposition}
\newtheorem{lemma}[theorem]{Lemma}
\newtheorem{corollary}[theorem]{Corollary}
\theoremstyle{definition}
\newtheorem{definition}[theorem]{Definition}
\theoremstyle{remark}
\newtheorem{remark}[theorem]{Remark}

\newcommand{\D}{\mathbb D}
\newcommand{\dd}{\,\mathrm d}
\newcommand{\sech}{\operatorname{sech}}
\newcommand{\csch}{\operatorname{csch}}
\newcommand{\Dom}{\operatorname{Dom}}
\newcommand{\Ran}{\operatorname{Ran}}
\newcommand{\spec}{\operatorname{spec}}
\newcommand{\norm}[1]{\lVert #1\rVert}
\newcommand{\ip}[2]{\langle #1,#2\rangle}
\newcommand{\todo}[1]{{\color{red!75!black}\textbf{[TODO:} #1\textbf{]}}}
\newcommand{\R}{\mathbb R}
\newcommand{\N}{\mathbb N}
\newcommand{\Z}{\mathbb Z}
\newcommand{\question}[1]{%
	\par\medskip\noindent
	\fbox{\parbox{0.94\linewidth}{\textbf{Question:} #1}}
	\par\medskip}
\numberwithin{equation}{section}

\hypersetup{
  pdftitle={Explicit diagonalization of the linearized Calderón problem on the disk: multiplicity normalization and exponential spectral attenuation},
  pdfauthor={William R. B. Lionheart},
  colorlinks=true,
  linkcolor=blue!45!black,
  citecolor=blue!45!black,
  urlcolor=blue!45!black
}

\begin{document}

\articletype{Paper}

\title{Explicit diagonalization of the linearized Calderón problem on the disk: multiplicity normalization and exponential spectral attenuation}

\author{William R B Lionheart$^{1,*}$}

\affil{$^1$Department of Mathematics, The University of Manchester, Manchester M13 9PL, United Kingdom}
\affil{$^*$Author to whom correspondence should be addressed.}
\email{bill.lionheart@manchester.ac.uk}

\keywords{Calder\'on problem, electrical impedance tomography, diagonalization,  Berezin transform, hyperbolic Laplacian, Zernike lattice, generalized singular functions}

\begin{abstract}
	
	We study the full angular dependence of the linearized Schr\"odinger
	form of the Calder\'on problem on the unit disk.  We obtain an explicit
	continuous spectral diagonalization of its normal operator and hence a
	generalized singular system for the linearized measurement map.
 After removing repeated boundary Fourier measurements, the independent data are indexed by the Zernike lattice.  Their normal operator is
\[
 (Kf)(z)=\frac1\pi\int_{\D}\frac{f(w)}{|1-z\overline w|^2}\,\dd A(w).
\]
A weighted Schur estimate gives the exact norm $\|K\|=\pi$, and a weakly null sequence of boundary annuli proves non-compactness before spectral theory is invoked.  We identify a second-order differential expression on the disk satisfying a direct global kernel commutator identity.  Multiplication by $(1-|z|^2)/2$ conjugates its canonical realization to the Poincar\'e Laplacian and conjugates $K$ to hyperbolic convolution with $(4\pi)^{-1}\sech^2(d_H/2)$.  This is a renormalized Hardy-endpoint member of the disk Berezin family.  Classical spherical harmonic analysis then yields
\[
 K=\pi\sech\!\left(\pi\sqrt{-\mathsf L-\tfrac14}\right),
\]
where $\mathsf L$ is the Euclidean-space realization unitarily equivalent to the hyperbolic Laplacian.  Hence $K$ has purely absolutely continuous spectrum $[0,\pi]$, and its generalized singular attenuation is exponential in hyperbolic frequency.  We give the associated Legendre and hypergeometric modes and a brief finite-section numerical check.  The classical Berezin diagonalization is not claimed as new; the main point is that this Hardy-endpoint operator arises naturally from multiplicity-normalized linearized Calder\'on measurements.
\end{abstract}

\noindent\todo{Confirm the final author list, affiliations, corresponding-author details and CRediT contributions. Done}

\section{Introduction}

Calder\'on introduced the boundary linearization that now bears his name in the conductivity setting \cite{Calderon1980}.  This is important for medical, industrial and geophysical imaging \cite{adler2015electrical}.

Numerical singular value decompositions of discretized versions of
this problem have been considered since 
 the 1990s \cite{breckon1990image}.  The radial subproblem, corresponding to the $m=0$ angular Fourier
block, is particularly transparent. For rotationally symmetric perturbations, a multiplicity-normalized normal operator reduces, after the change of variables $t=r^2$, to the Hilbert kernel \cite{Lionheart2026AtLast}. 
\begin{equation}
 (Hf)(t)=\int_0^1\frac{f(s)}{1-st}\,\dd s.
 \label{eq:radial-hilbert}
\end{equation}
Its continuous spectral multiplier may be obtained from an adjacent-interval finite Hilbert transform, from a commuting Sturm--Liouville expression, from Hilbert-matrix theory, or from a P\"oschl--Teller scattering problem; see, for example, \cite{KatsevichTovbis2016,AlemanMontesSarafoleanu2012,MontesVirtanen2025}.  The purpose of the present paper is to show that the full, non-rotationally symmetric problem is not merely a direct sum of unrelated radial calculations.  The angular blocks assemble into one hyperbolically invariant operator on the disk and this gives a diagonalization of this non-compact operator. 

The key bookkeeping observation is that a matrix element indexed by boundary Fourier frequencies $(k,l)$ depends on the two integers
\begin{equation}
 m=k-l,\qquad n=|k|+|l|.
 \label{eq:intro-invariants}
\end{equation}
The admissible pairs satisfy
\[
 n\geq |m|,\qquad n\equiv m\pmod 2,
\]
which is exactly the Zernike index lattice.  Counting every independent datum once, rather than with its raw matrix multiplicity, gives the moment map whose normal kernel is
\begin{equation}
 \kappa(z,w)=\frac1{\pi|1-z\overline w|^2}.
 \label{eq:intro-kernel}
\end{equation}
This kernel is not the Bergman projection kernel.  After a unitary change of measure it is instead the squared overlap of normalized Hardy reproducing kernels and is a limiting Berezin convolution.  Berezin transforms and their expression as functions of invariant differential operators have a substantial literature \cite{Berezin1975,UnterbergerUpmeier1994,Englis1995,Stroethoff1997}.  Our contribution is the bridge from linearized Calder\'on measurements to the Hardy-endpoint kernel, together with its inverse-problem interpretation.

The main results may be summarized as follows.
\begin{enumerate}[label=(\roman*)]
\item The Schr\"odinger Dirichlet-to-Neumann map is Fr\'echet differentiable at the zero potential for perturbations in $L^2(\D)$, and its derivative is the usual Alessandrini bilinear form.
\item The multiplicity-reduced Fourier data define a bounded moment map $\mathcal{M}$ with 
\[
 K=\mathcal{M}^*\mathcal{M},
 \]
and kernel \eqref{eq:intro-kernel}.
\item The operator is positive and self-adjoint, has exact norm $\pi$, and is non-compact.
\item A global second-order differential expression $\mathcal P$ satisfies
\[
 \mathcal P_z\kappa(z,w)=\mathcal P_w\kappa(z,w),
\]
not merely after separation into angular modes.
\item Under a unitary multiplication operator $U$, the auxiliary expression $\mathcal L=\mathcal P-\frac12\partial_\theta^2-1$ becomes the Poincar\'e Laplacian and $UKU^{-1}$ becomes a radial hyperbolic convolution.
\item The exact normal-operator multiplier is
\[
\lambda(\mu)=\frac{\pi}{\cosh(\pi\mu)},
\]
so that the spectrum is continuous and high hyperbolic frequencies
are exponentially attenuated.
\item The generalized eigenfunctions of $K$, and hence the generalized
right singular functions of $\mathcal M$, are given explicitly by
associated Legendre or hypergeometric functions.  Their Liouville
normal form is a P\"oschl--Teller scattering equation, which explains
the concentration of fine radial oscillations near the boundary.

\end{enumerate}

The distinction between multiplicity normalization and physical frequency weighting is important.  The former only removes repeated copies of identical linearized measurements.  A Sobolev, noise-whitening, Dirichlet-to-Neumann, or Neumann-to-Dirichlet weighting may introduce additional factors depending on $k$ and $l$ and therefore changes the normal operator.  The paper treats the equal-weight, non-redundant data space exactly.  \todo{After checking the original EIT normalization, state explicitly how this data norm relates to the physical DtN and NtD norms.}

\section{Linearization in an $L^2$ potential}
\label{sec:linearization}

Let $\D\subset\R^2$ be the unit disk and consider
\begin{equation}
 (-\Delta+q)u=0\quad\hbox{in }\D,
 \qquad u|_{\partial\D}=f.
 \label{eq:schrodinger-pde}
\end{equation}
For $q$ close to zero in $L^2(\D)$, write $\Lambda_q:H^{1/2}(\partial\D)\to H^{-1/2}(\partial\D)$ for the Dirichlet-to-Neumann map, with the outward normal convention.  Note that while this Schr\"odinger version has uses, the traditional problem is framed in terms of a conductivity, related to $\eta$ by a Liouville transformation. 
For a positive conductivity $\gamma$, the substitution
$u=\gamma^{1/2}v$ transforms
\[
\nabla\cdot(\gamma\nabla v)=0
\]
into
\[
(-\Delta+q_\gamma)u=0,
\qquad
q_\gamma=\gamma^{-1/2}\Delta\gamma^{1/2}.
\]
If $\gamma=1$ in a neighbourhood of $\partial\D$, the corresponding
boundary maps agree under this transformation.  For general boundary
values of $\gamma$ there is also the standard boundary conjugation and
zeroth-order correction.

\begin{proposition}[Linearization at the zero potential]
\label{prop:linearization}
There is a neighbourhood of zero in $L^2(\D)$ on which \eqref{eq:schrodinger-pde} is uniquely solvable and $q\mapsto\Lambda_q$ is analytic as a map
\[
 L^2(\D)\longrightarrow
 \mathcal L\bigl(H^{1/2}(\partial\D),H^{-1/2}(\partial\D)\bigr).
\]
In particular,
\begin{equation}
 \ip{D\Lambda_0[\eta]f}{g}
 =\int_{\D}\eta\,u_f^0\,\overline{u_g^0}\,\dd A,
 \label{eq:linearized-DtN}
\end{equation}
where $u_f^0$ and $u_g^0$ are harmonic extensions.  Moreover,
\begin{equation}
 \norm{D\Lambda_0[\eta]}_{H^{1/2}\to H^{-1/2}}
 \leq C\norm{\eta}_{L^2},
 \label{eq:derivative-bound}
\end{equation}
and
\begin{equation}
 \norm{\Lambda_q-\Lambda_0-D\Lambda_0[q]}_{H^{1/2}\to H^{-1/2}}
 \leq C\norm{q}_{L^2}^2
 \label{eq:quadratic-remainder}
\end{equation}
for $\norm q_{L^2}$ sufficiently small.
\end{proposition}

\begin{proof}
Multiplication by $q$ defines a bounded map $M_q:H^1(\D)\to H^{-1}(\D)$.  Indeed, the two-dimensional embedding $H^1(\D)\hookrightarrow L^4(\D)$ and Cauchy--Schwarz give
\begin{equation}
 \left|\int_{\D}q u\overline v\,\dd A\right|
 \leq \norm q_2\norm u_4\norm v_4
 \leq C\norm q_2\norm u_{H^1}\norm v_{H^1}.
 \label{eq:multiplication-estimate}
\end{equation}
Let $G=(-\Delta_D)^{-1}:H^{-1}(\D)\to H_0^1(\D)$.  If $u_f^0$ denotes harmonic extension, the solution is
\begin{equation}
 u_f^q=(I+GM_q)^{-1}u_f^0.
 \label{eq:operator-series}
\end{equation}
For small $\norm q_2$, the inverse is represented by a norm-convergent Neumann series, proving analyticity and
\[
 \norm{u_f^q-u_f^0}_{H^1}
 \leq C\norm q_2\norm f_{H^{1/2}}.
\]
The integration-by-parts formula commonly called Alessandrini's identity is
\begin{equation}
 \ip{(\Lambda_q-\Lambda_0)f}{g}
 =\int_{\D}q\,u_f^q\,\overline{u_g^0}\,\dd A.
 \label{eq:alessandrini}
\end{equation}
Subtraction of \eqref{eq:linearized-DtN}, followed by \eqref{eq:multiplication-estimate}, gives \eqref{eq:quadratic-remainder}.  The same estimate with harmonic extensions gives \eqref{eq:derivative-bound}.
\end{proof}

\begin{remark}
The operator-series argument gives every Fr\'echet derivative, but only the first derivative is needed below.  We use \eqref{eq:alessandrini} in the standard descriptive sense; it is, at heart, the Green identity for two solutions \cite{Alessandrini1988}.
\end{remark}

Let
\begin{equation}
 \varphi_k(\theta)=(2\pi)^{-1/2}e^{ik\theta},\qquad k\in\Z.
 \label{eq:boundary-fourier}
\end{equation}
Its harmonic extension is $u_k^0(r,\theta)=r^{|k|}\varphi_k(\theta)$.  If
\[
 \eta_m(r)=\frac1{2\pi}\int_0^{2\pi}\eta(r,\theta)e^{-im\theta}\,\dd\theta,
\]
then \eqref{eq:linearized-DtN} yields
\begin{equation}
 \ip{D\Lambda_0[\eta]\varphi_k}{\varphi_l}
 =\int_0^1\eta_{l-k}(r)r^{|k|+|l|}r\,\dd r.
 \label{eq:fourier-matrix-element}
\end{equation}
Thus, up to the harmless sign convention for the object angular index, the datum depends only on $k-l$ and $|k|+|l|$.

\todo{Check equation \eqref{eq:fourier-matrix-element} against the exact Fourier, complex-conjugation and DtN sign conventions used in the intended submission, including any conductivity-to-Schr\"odinger normalization.}

\section{Non-redundant measurements and the Zernike lattice}
\label{sec:nonredundant}

Define
\begin{equation}
 m=k-l,\qquad n=|k|+|l|.
 \label{eq:mn}
\end{equation}
For fixed $(m,n)$ let
\[
 F_{m,n}=\{(k,l)\in\Z^2:k-l=m,\ |k|+|l|=n\}.
\]

\begin{proposition}[Fibre classification]
\label{prop:fibres}
The fibre is nonempty precisely when
\begin{equation}
 n\geq |m|,\qquad n\equiv m\pmod 2.
 \label{eq:zernike-lattice}
\end{equation}
For admissible $(m,n)$,
\begin{equation}
 \#F_{m,n}=\nu(m,n)=
 \begin{cases}
 |m|+1,&n=|m|,\\
 2,&n>|m|.
 \end{cases}
 \label{eq:fibre-multiplicity}
\end{equation}
\end{proposition}

\begin{proof}
Put $v=k+l$.  The identity
\begin{equation}
 |k|+|l|=\max\{|k-l|,|k+l|\}
 \label{eq:max-identity}
\end{equation}
gives $n=\max\{|m|,|v|\}$.  Moreover,
\[
 k=\frac{m+v}{2},\qquad l=\frac{v-m}{2},
\]
so $v$ and $m$ have the same parity.  If $n>|m|$, then $v=\pm n$, giving two points.  If $n=|m|$, then $v=-n,-n+2,\ldots,n$, giving $n+1=|m|+1$ points.
\end{proof}

The sum-and-difference coordinates
\begin{equation}
 p=\frac{n+m}{2},\qquad q=\frac{n-m}{2}
 \label{eq:pq}
\end{equation}
give a bijection between the lattice \eqref{eq:zernike-lattice} and $\N_0^2$, with
\[
 n=p+q,\qquad m=p-q,
\]
and radial Zernike index
\begin{equation}
 j=\min\{p,q\}=\frac{n-|m|}{2}.
 \label{eq:radial-index}
\end{equation}
This is the familiar bidegree-to-polar-degree relation
\[
 z^p\overline z^{\,q}=r^{p+q}e^{i(p-q)\theta}.
\]

When a redundant data matrix is constant on each fibre, the norm
\begin{equation}
 \norm Y_\sharp^2=
 \sum_{k,l\in\Z}\frac{|Y_{kl}|^2}{\nu(k-l,|k|+|l|)}
 \label{eq:sharp-data-norm}
\end{equation}
counts every independent datum exactly once.  Equivalently, one may select one representative for each $(m,n)$ or $(p,q)$.

Define the non-redundant moment map initially on a dense subspace by
\begin{equation}
 \mathcal{M} f_{p,q}
 =\frac1{\sqrt\pi}\int_{\D}
 f(w)\overline w^{\,p}w^q\,\dd A(w),
 \qquad p,q\in\N_0.
 \label{eq:M}
\end{equation}
In $(m,j)$ coordinates this is
\begin{equation}
 (\mathcal{M}  f)_{m,j}
 =\frac1{\sqrt\pi}\int_{\D}
 f(re^{i\theta})r^{|m|+2j}e^{-im\theta}\,\dd A,
 \qquad m\in\Z,\quad j\in\N_0.
 \label{eq:sigma-mj}
\end{equation}
The normalization in \eqref{eq:M} is chosen to make the later kernel and norm formulas simple.


\section{The normal operator}
\label{sec:normal}

\begin{proposition}[Normal-operator factorization]
\label{prop:normal-factorization}
The map $\mathcal{M} $ extends to a bounded operator
\[
 \norm{\mathcal M}=\sqrt{\pi}:L^2(\D,\dd A)\longrightarrow\ell^2(\N_0^2),
\]
and
\begin{equation}
 K=\mathcal{M} ^*\mathcal{M},
 \qquad
 (Kf)(z)=\frac1\pi\int_{\D}
 \frac{f(w)}{|1-z\overline w|^2}\,\dd A(w).
 \label{eq:K-definition}
\end{equation}
In angular frequency $m$, the normal kernel is
\begin{equation}
 k_m(r,\rho)=\frac{(r\rho)^{|m|}}{1-r^2\rho^2}.
 \label{eq:angular-block-r}
\end{equation}
\end{proposition}

\begin{proof}
For a finitely supported sequence $c=(c_{p,q})$,
\[
 ( \mathcal{M}^*c)(z)
 =\frac1{\sqrt\pi}\sum_{p,q\geq0}c_{p,q}z^p\overline z^{\,q}.
\]
Consequently, finite partial sums give the kernel
\[
 \frac1\pi\sum_{p,q\geq0}z^p\overline z^{\,q}
 \overline w^{\,p}w^q
 =\frac1{\pi|1-z\overline w|^2}.
\]
The boundedness proved in Proposition~\ref{prop:norm} then extends the quadratic-form identity to all of $L^2$.  For fixed $m$, summing over $j$ gives
\[
 \sum_{j=0}^\infty(r\rho)^{|m|+2j}
 =\frac{(r\rho)^{|m|}}{1-r^2\rho^2}.
\]
\end{proof}

The Poisson-kernel identity reassembles the angular blocks into the
global kernel:
\question{Not really reconstruction?}
\begin{equation}
 \sum_{m\in\Z}\frac{(r\rho)^{|m|}}{1-r^2\rho^2}
 e^{im(\theta-\phi)}
 =\frac1{|1-z\overline w|^2}.
 \label{eq:block-sum}
\end{equation}
Under $t=r^2$ and $s=\rho^2$, the $m$th radial block is unitarily equivalent to
\begin{equation}
 (K_mh)(t)=\int_0^1\frac{(st)^{|m|/2}}{1-st}h(s)\,\dd s.
 \label{eq:block-t}
\end{equation}
The case $m=0$ is exactly \eqref{eq:radial-hilbert}.

\begin{remark}[Raw multiplicity]
If the unreduced matrix is given the ordinary entrywise $\ell^2$ norm, the edge $n=|m|$ is over-counted.  In each angular block the raw normal kernel differs from a constant multiple of \eqref{eq:angular-block-r} by a rank-one term proportional to $r^{|m|}\rho^{|m|}$.  The normalization \eqref{eq:sharp-data-norm} removes precisely these exceptional edge terms.  \todo{Write the exact formula using the final Fourier normalization and decide whether it belongs in the main text or an appendix.}
\end{remark}

\section{Exact norm and non-compactness}
\label{sec:elementary}

The following arguments establish the operator before any spectral diagonalization is used.

\begin{proposition}[Weighted Schur test]
\label{prop:norm}
The operator $K$ is bounded, positive and self-adjoint on $L^2(\D,\dd A)$, and
\begin{equation}
 \norm K=\pi.
 \label{eq:exact-norm}
\end{equation}
Consequently, $\norm{\mathcal{M}}=\sqrt\pi$.
\end{proposition}

\begin{proof}
For the non-negative symmetric kernel
\[
 \kappa(z,w)=\frac1{\pi|1-z\overline w|^2},
\]
we use the Schur test  in the form given, for example, in
\cite[Section~4]{HalmosSunder1978} with  weight $\chi(z)=(1-|z|^2)^{-1/2}$.  With $z=re^{i\theta}$ and $w=\rho e^{i\phi}$,
\[
 \int_0^{2\pi}\frac{\dd\phi}
 {1-2r\rho\cos\phi+r^2\rho^2}
 =\frac{2\pi}{1-r^2\rho^2}.
\]
Therefore
\begin{align}
 \int_{\D}\kappa(z,w)\chi(w)\,\dd A(w)
 &=2\int_0^1\frac{\rho\,\dd\rho}
 {(1-r^2\rho^2)\sqrt{1-\rho^2}}\notag\\
 &=\frac{2\arcsin r}{r\sqrt{1-r^2}}
 \leq \pi\chi(z).
 \label{eq:schur-calculation}
\end{align}
Schur's test gives $\norm K\leq\pi$.

The closed radial subspace reduces $K$.  The unitary map
\[
 (Wf)(t)=\sqrt\pi f(\sqrt t)
\]
identifies the radial restriction with $H$ in \eqref{eq:radial-hilbert}.  The standard near-extremizers $(1-t)^{-1/2+\epsilon}$ show $\norm H\geq\pi$, and hence $\norm K\geq\pi$.  Positivity follows from the factorization in Proposition~\ref{prop:normal-factorization}, and self-adjointness follows from symmetry and boundedness.
\end{proof}

\begin{proposition}[Boundary-annulus non-compactness]
\label{prop:noncompact}
The operator $K$ is not compact.
\end{proposition}

\begin{proof}
For $0<a<1$ set
\begin{equation}
 F_a(z)=\frac1{\sqrt{\pi a}}
 \mathbf1_{\{1-a<|z|^2<1\}}(z).
 \label{eq:annulus-sequence}
\end{equation}
Then $\norm{F_a}_2=1$ and $F_a\rightharpoonup0$ as $a\downarrow0$.  Under the radial unitary map, $F_a$ becomes
\[
 h_a(t)=a^{-1/2}\mathbf1_{(1-a,1)}(t).
\]
A direct calculation gives
\[
 (Hh_a)(t)=\frac1{\sqrt a\,t}
 \log\!\left(\frac{1-(1-a)t}{1-t}\right).
\]
Using $x=at/(1-t)$,
\begin{align*}
 \norm{Hh_a}_2^2
 &=\int_0^\infty\frac{\log^2(1+x)}{x^2}\,\dd x\\
 &=2\int_0^\infty\frac{\log(1+x)}{x(1+x)}\,\dd x
 =2\zeta(2)=\frac{\pi^2}{3}.
\end{align*}
Thus $\norm{KF_a}=\pi/\sqrt3$ for every $a$.  A compact operator maps a bounded weakly null sequence to a norm-null sequence, which is impossible here.
\end{proof}

\section{A global differential commutator}
\label{sec:commutator}

Set $t=r^2$.  The angular blocks \eqref{eq:block-t} suggest the differential expression
\begin{equation}
 \mathcal P
 =\partial_t\!\left(t(1-t)^2\partial_t\right)+t
 +\frac14\left(t+\frac1t\right)\partial_\theta^2.
 \label{eq:P-t}
\end{equation}
In $(r,\theta)$ variables,
\begin{equation}
 4\mathcal P
 =\frac1r\partial_r\!\left(r(1-r^2)^2\partial_r\right)
 +\left(r^2+\frac1{r^2}\right)\partial_\theta^2+4r^2.
 \label{eq:P-r}
\end{equation}
This expression is formally self-adjoint with respect to Euclidean area $r\,\dd r\,\dd\theta$.

\begin{proposition}[Direct kernel identity]
\label{prop:kernel-commutator}
For $z,w\in\D$,
\begin{equation}
 \mathcal P_z\kappa(z,w)=\mathcal P_w\kappa(z,w),
 \qquad
 \kappa(z,w)=\frac1{\pi|1-z\overline w|^2}.
 \label{eq:kernel-commutator}
\end{equation}
Consequently,
\begin{equation}
 \mathcal P Kf=K\mathcal P f,
 \qquad f\in C_c^\infty(\D).
 \label{eq:core-commutation}
\end{equation}
\end{proposition}

\begin{proof}
Let $\omega(z)=1-|z|^2$ and define
\begin{equation}
 \mathcal L=\mathcal P-\frac12\partial_\theta^2-1.
 \label{eq:L-def}
\end{equation}
A direct expansion gives
\begin{equation}
 \mathcal Lf=\frac{\omega}{4}\Delta_E(\omega f).
 \label{eq:L-euclidean}
\end{equation}
Introduce
\begin{equation}
 \delta(z,w)=\frac{\omega(z)\omega(w)}{|1-z\overline w|^2}.
 \label{eq:delta}
\end{equation}
For the Poincar\'e metric introduced in Section~\ref{sec:hyperbolic}, $\delta(z,w)=\sech^2(d_H(z,w)/2)$.  If $F(d)=\sech^2(d/2)$, then
\begin{equation}
 F''(d)+\coth(d)F'(d)=-F(d)^2.
 \label{eq:radial-delta-identity}
\end{equation}
Hence $\Delta_{H,z}\delta=-\delta^2=\Delta_{H,w}\delta$.  Using \eqref{eq:L-euclidean}, this gives
\begin{equation}
 \mathcal L_z\kappa(z,w)
 =-\frac{\omega(z)\omega(w)}{\pi|1-z\overline w|^4}
 =\mathcal L_w\kappa(z,w).
 \label{eq:L-kernel}
\end{equation}
The kernel depends on the angular variables only through $\theta-\phi$, so $\partial_\theta^2\kappa=\partial_\phi^2\kappa$.  Equations \eqref{eq:L-def} and \eqref{eq:L-kernel} prove \eqref{eq:kernel-commutator}.  For compactly supported $f$, formal self-adjointness and integration by parts have no boundary contribution, giving \eqref{eq:core-commutation}.
\end{proof}

\todo{Run the independent Mathematica check of \eqref{eq:kernel-commutator} in polar coordinates, retain the script with the paper archive, and compare its conventions with \eqref{eq:P-r}.}

Equation \eqref{eq:core-commutation} alone does not specify a self-adjoint realization of the unbounded differential expression.  The canonical domain and full operator commutation will follow from the hyperbolic realization.

\section{Hyperbolic and Berezin realization}
\label{sec:hyperbolic}

Equip $\D$ with the Poincar\'e metric of curvature $-1$,
\begin{equation}
 \dd s_H^2=\frac{4|\dd z|^2}{(1-|z|^2)^2},
 \qquad
 \dd\mu_H=\frac{4\,\dd A}{(1-|z|^2)^2},
 \label{eq:poincare-metric}
\end{equation}
and Laplace--Beltrami operator
\begin{equation}
 \Delta_H=\frac{(1-|z|^2)^2}{4}\Delta_E.
 \label{eq:hyperbolic-laplacian}
\end{equation}
Our sign convention makes $\Delta_H$ non-positive on $L^2(\D,\dd\mu_H)$.

Define
\begin{equation}
 U:L^2(\D,\dd A)\longrightarrow L^2(\D,\dd\mu_H),
 \qquad
 (Uf)(z)=\frac{1-|z|^2}{2}f(z).
 \label{eq:U}
\end{equation}
This map is unitary.

\begin{theorem}[Hyperbolic conjugation]
\label{thm:hyperbolic-conjugation}
Let $\mathsf L$ be the self-adjoint realization of \eqref{eq:L-euclidean} defined by
\begin{equation}
 \mathsf L=U^{-1}\Delta_HU.
 \label{eq:L-operator}
\end{equation}
Then
\begin{equation}
 (UKU^{-1}g)(z)
 =\frac1{4\pi}\int_{\D}
 \sech^2\!\left(\frac{d_H(z,w)}2\right)
 g(w)\,\dd\mu_H(w).
 \label{eq:hyperbolic-convolution}
\end{equation}
Thus $UKU^{-1}$ is an isometry-invariant radial convolution on the hyperbolic plane.
\end{theorem}

\begin{proof}
The first assertion is \eqref{eq:L-euclidean} and the definitions of $U$ and $\Delta_H$.  The hyperbolic distance satisfies
\begin{equation}
 \cosh^2\!\left(\frac{d_H(z,w)}2\right)
 =\frac{|1-z\overline w|^2}
 {(1-|z|^2)(1-|w|^2)}.
 \label{eq:distance-identity}
\end{equation}
Substitution of \eqref{eq:U} and \eqref{eq:distance-identity} into \eqref{eq:K-definition} gives \eqref{eq:hyperbolic-convolution}.
\end{proof}

The invariant factor
\begin{equation}
 \delta(z,w)=\frac{(1-|z|^2)(1-|w|^2)}{|1-z\overline w|^2}
 \label{eq:hardy-overlap}
\end{equation}
is the squared overlap of normalized Hardy kernels
\[
 s_z(\zeta)=\frac{\sqrt{1-|z|^2}}{1-\overline z\zeta},
 \qquad
 |\ip{s_z}{s_w}_{H^2}|^2=\delta(z,w).
\]
For $\lambda>1$, the weighted Bergman Berezin transform has invariant kernel proportional to $\delta^\lambda$.  Formula \eqref{eq:hyperbolic-convolution} is the renormalized endpoint $\lambda\downarrow1$.  We shall therefore call it the Hardy-endpoint Berezin convolution.  This terminology distinguishes it from the Bergman projection, which is an idempotent holomorphic projection and has spectrum $\{0,1\}$; the Berezin transform is an averaging operator and is not a projection \cite{Stroethoff1997,Englis1995}.

The endpoint kernel
\[
 h(d)=\frac1{4\pi}\sech^2(d/2)
\]
is not in $L^1(\mathbb H^2)$: $2\pi h(d)\sinh d\to1$ as $d\to\infty$.  It is in $L^2$.  Thus boundedness should not be justified by an $L^1$ Young inequality.  Proposition~\ref{prop:norm} establishes boundedness independently, and the spherical transform below identifies the bounded multiplier.

\todo{Locate and cite the most explicit classical source for the precise $\lambda\downarrow1$ endpoint and its normalization; search the Berezin/representation-theory literature, including Russian-language antecedents, before making any priority statement about the transform formula itself.}

\section{Spherical multiplier and operator spectrum}
\label{sec:spectrum}

The spectrum of $-\Delta_H$ on the hyperbolic plane is $[1/4,\infty)$ and is purely absolutely continuous \cite{Helgason1984}.  Write
\begin{equation}
 \mathsf A=-\mathsf L-\frac14I\geq0.
 \label{eq:A}
\end{equation}
The spherical eigenfunction with parameter $\mu\geq0$, normalized to be one at the origin, is
\[
 \phi_\mu(d)=P_{-1/2+i\mu}(\cosh d),
 \qquad
 \Delta_H\phi_\mu=-\left(\mu^2+\frac14\right)\phi_\mu.
\]

\begin{lemma}[Endpoint spherical transform]
\label{lem:spherical-transform}
For $\mu\geq0$,
\begin{equation}
 \pi\int_1^\infty\frac{P_{-1/2+i\mu}(u)}{1+u}\,\dd u
 =\frac{\pi^2}{\cosh(\pi\mu)}.
 \label{eq:mehler-fock-integral}
\end{equation}
With the normalization in \eqref{eq:hyperbolic-convolution}, the corresponding multiplier is
\begin{equation}
 \widehat h(\mu)=\frac{\pi}{\cosh(\pi\mu)}.
 \label{eq:multiplier}
\end{equation}
\end{lemma}

\begin{proof}
In polar coordinates about $z$, the spherical transform of \eqref{eq:hyperbolic-convolution} is
\[
 2\pi\int_0^\infty\frac1{4\pi}\sech^2(d/2)
 P_{-1/2+i\mu}(\cosh d)\sinh d\,\dd d.
\]
Since $\frac12\sech^2(d/2)\sinh d=\frac{2\sinh d}{2(1+\cosh d)}=\frac{\sinh d}{1+\cosh d}$, the substitution $u=\cosh d$ reduces this to
\[
 \int_1^\infty\frac{P_{-1/2+i\mu}(u)}{1+u}\,\dd u.
\]
The Mehler--Fock integral is the first equality in \eqref{eq:mehler-fock-integral}; its value follows from the hypergeometric representation and Euler's beta integral, or from the standard rank-one Berezin multiplier formula.  The factor $1/\pi$ in \eqref{eq:K-definition} gives \eqref{eq:multiplier}.
\end{proof}

\todo{Check every factor of $2$ and $\pi$ in Lemma~\ref{lem:spherical-transform} against a cited Mehler--Fock or Berezin formula.  The independent checks $\widehat h(0)=\pi$ and $\norm K=\pi$ must agree.}

\begin{theorem}[Exact spectral formula]
\label{thm:spectral-formula}
On $L^2(\D,\dd A)$,
\begin{equation}
 \boxed{\quad
 K=\pi\sech\!\left(\pi\sqrt{-\mathsf L-\tfrac14I}\right).
 \quad}
 \label{eq:exact-spectral-formula}
\end{equation}
In particular,
\begin{equation}
 \spec(K)=[0,\pi],
 \label{eq:spectrum-K}
\end{equation}
$K$ has purely absolutely continuous spectrum, is injective, and has dense non-closed range.  Neither $0$ nor $\pi$ is an $L^2$ eigenvalue.
\end{theorem}

\begin{proof}
The radial convolution in Theorem~\ref{thm:hyperbolic-conjugation} is diagonalized by the spherical/Helgason Fourier transform.  Lemma~\ref{lem:spherical-transform} gives the bounded multiplier $\pi\sech(\pi\mu)$, yielding \eqref{eq:exact-spectral-formula}.  The multiplier is continuous, strictly positive for finite $\mu$, decreases from $\pi$ to zero, and the hyperbolic Laplacian has purely absolutely continuous spectrum.  The stated consequences follow from the spectral theorem.
\end{proof}

The formula also supplies the canonical closed realization of the commutator.  Let
\begin{equation}
 J=-i\partial_\theta.
 \label{eq:J}
\end{equation}
Rotations are hyperbolic isometries, so $\mathsf A$ and $J$ strongly commute.  Define
\begin{equation}
 \mathsf P=\frac34I-\mathsf A-\frac12J^2,
 \qquad
 \Dom(\mathsf P)=\Dom(\mathsf A)\cap\Dom(J^2).
 \label{eq:P-operator}
\end{equation}

\begin{corollary}[Strong commutation]
\label{cor:strong-commutation}
The operator $\mathsf P$ is self-adjoint, agrees with \eqref{eq:P-t} on $C_c^\infty(\D)$, and
\begin{equation}
 K\Dom(\mathsf P)\subseteq\Dom(\mathsf P),
 \qquad
 \mathsf P Kf=K\mathsf P f
 \quad(f\in\Dom(\mathsf P)).
 \label{eq:strong-commutation}
\end{equation}
Indeed, $K$ strongly commutes with the spectral projections of $\mathsf P$.
\end{corollary}

\begin{proof}
In the joint spectral representation of $(\mathsf A,J)$, the operators $K$ and $\mathsf P$ are multiplication by
\[
 \frac{\pi}{\cosh(\pi\mu)}
 \quad\hbox{and}\quad
 \frac34-\mu^2-\frac{m^2}{2},
\]
respectively.  The first multiplier is bounded, so it preserves the domain defined by square integrability against the second multiplier and commutes pointwise with it.
\end{proof}

\section{Generalized singular functions and resolution}
\label{sec:singular-functions}

Let $x=d_H(0,r)=2\operatorname{artanh}r$, so $r=\tanh(x/2)$.  Separating angular frequency $m\in\Z$, a regular generalized eigenfunction of the hyperbolic Laplacian is
\begin{align}
 R_{m,\mu}(x)
 &=\tanh(x/2)^{|m|}
 {}_2F_1\!\left(\frac12+i\mu,\frac12-i\mu;
 |m|+1;-\sinh^2(x/2)\right)\notag\\
 &=\Gamma(|m|+1)P_{-1/2+i\mu}^{-|m|}(\cosh x).
 \label{eq:radial-genfun}
\end{align}
The corresponding Euclidean-space generalized function is
\begin{equation}
 \Phi_{m,\mu}(r,\theta)
 =\frac{2}{1-r^2}R_{m,\mu}(x)e^{im\theta}.
 \label{eq:euclidean-genfun}
\end{equation}
Formally,
\begin{equation}
 K\Phi_{m,\mu}
 =\frac{\pi}{\cosh(\pi\mu)}\Phi_{m,\mu}.
 \label{eq:gen-eigen-K}
\end{equation}
These are generalized, not square-integrable, eigenfunctions.

Because $K=(\Sigma^\sharp)^*\Sigma^\sharp$, the generalized singular attenuation of the measurement map is
\begin{equation}
 \sigma(\mu)=\left(\frac{\pi}{\cosh(\pi\mu)}\right)^{1/2}
 \sim\sqrt{2\pi}\,e^{-\pi\mu/2}
 \qquad(\mu\to\infty).
 \label{eq:singular-attenuation}
\end{equation}
Thus an $L^2$ perturbation component at hyperbolic frequency $\mu$ is exponentially suppressed in the data.  In a deterministic noise model $d^\delta=\Sigma^\sharp f+e$, $\norm e\leq\delta$, inversion of that component amplifies noise by approximately $\sigma(\mu)^{-1}$.  This is the appropriate continuous-spectrum analogue of exponentially decaying singular values.

The stable Liouville profile
\begin{equation}
 v_{m,\mu}(x)=\sqrt{\sinh x}\,R_{m,\mu}(x)
 \label{eq:liouville-profile}
\end{equation}
satisfies
\begin{equation}
 -v''+\left(m^2-\frac14\right)\csch^2x\,v=\mu^2v.
 \label{eq:poschl-teller}
\end{equation}
The potential is short range as $x\to\infty$, so $v$ is asymptotically oscillatory with wavelength $\pi/\mu$ between successive zeros.  In Euclidean radius, equal increments of $x$ are compressed exponentially towards $r=1$.  This gives a direct geometric explanation of the concentration of fine resolution near the boundary.

The exponential multiplier should not be confused with compactness.  It is a nonzero multiplication function on a continuous spectral space; multiplication by such a function is not compact.  Proposition~\ref{prop:noncompact} gives an elementary manifestation of the same fact.

\section{A brief finite-section experiment}
\label{sec:numerics}

For illustration, we truncate the hyperbolic radial variable to
$0<x<X$.  For each angular frequency $|m|\leq M$, a Gauss--Legendre
quadrature rule is used to form a symmetric Nystr\"om matrix for the
corresponding truncated integral operator.  If $x_i$ and $w_i$ are
the quadrature nodes and weights and $\widetilde k_m$ denotes the
kernel in the transformed coordinate, the matrix has the form
\[
(A_m)_{ij}
=
\sqrt{w_i}\,\widetilde k_m(x_i,x_j)\sqrt{w_j}.
\]
The leading eigenpairs of the real symmetric matrix $A_m$ are computed
using MATLAB's \texttt{eigs} routine.  If
$\lambda_{m,j}^{(X,N)}$ is a computed eigenvalue, then
\[
\sigma_{m,j}^{(X,N)}
=
\sqrt{\lambda_{m,j}^{(X,N)}}
\]
is the corresponding singular value of the truncated measurement
map.


The truncation makes every block compact; its discrete modes should be interpreted as spectral packets sampling the continuous parameter $\mu$, not as convergent eigenvalues of the infinite operator.


Figure~\ref{fig:global-order} shows the characteristic interlacing of angular frequency and radial oscillation count.  The sawtooth behaviour is expected: after ordering by attenuation, the next mode need not have either the next angular frequency or the next radial zero count.  Figure~\ref{fig:analytic-comparison} compares selected numerical profiles with \eqref{eq:radial-genfun} after the stable transformation \eqref{eq:liouville-profile}.

\begin{figure}[htbp]
 \centering
 \includegraphics[width=0.93\textwidth]{figures/nonradial_global_singular_values.png}
 \caption{Finite-section global ordering.  The singular values are compared with the analytic multiplier, while the lower panels display the angular frequency and radial-zero count of the correspondingly ordered mode.  The discrete ordering depends on the hyperbolic truncation, but the mixed angular--radial pattern is robust.}
 \label{fig:global-order}
\end{figure}

\begin{figure}[htbp]
 \centering
 \includegraphics[width=0.93\textwidth]{figures/nonradial_analytic_comparison.png}
 \caption{Comparison of finite-section numerical modes and associated Legendre/hypergeometric profiles in the Liouville variable.  The agreement away from the artificial boundary is a stronger check than singular-value matching alone.}
 \label{fig:analytic-comparison}
\end{figure}
\FloatBarrier

\todo{Repeat the convergence study under independent changes of $X$, radial quadrature size and angular cutoff; quantify residuals and profile errors; document the treatment of near-degenerate modes.  The present figures are exploratory, not a numerical-analysis theorem.}

\section{Discussion}
\label{sec:discussion}

The result supplies a global counterpart of the rotationally symmetric calculation.  The one-dimensional Hilbert-kernel operator is the $m=0$ spherical sector of a Hardy-endpoint Berezin convolution on the Poincar\'e disk.  The separate Sturm--Liouville expressions are the angular reductions of one disk differential operator.  This explains why the same multiplier appears through finite Hilbert transforms, Hilbert matrices, scattering theory and hyperbolic harmonic analysis.

The relation with disk Radon transforms is structural rather than literal.  In Radon and geodesic X-ray problems, Zernike or generalized Zernike functions diagonalize distinguished disk operators, and differential intertwining identities yield sharp mapping results \cite{Monard2020,MishraMonard2021,MishraMonardZou2023,EptaminitakisMonardZou2026}.  Here the same Zernike lattice arises from the multiplicities of boundary Fourier measurements, but the normal operator is not an X-ray normal operator and does not localize the unknown along hyperbolic geodesics.  The geometry is hyperbolic while the imaging mechanism is an exponentially smoothing invariant convolution.

The abstract diagonalization belongs to established Berezin and symmetric-space harmonic analysis.  The new inverse-problem statement is the route
\begin{equation*}
\begin{gathered}
 \hbox{linearized Calder\'on matrix}
 \longrightarrow \hbox{multiplicity fibres}
 \longrightarrow \hbox{Zernike lattice}\\
 \longrightarrow \frac1{\pi|1-z\overline w|^2}
 \longrightarrow \hbox{Hardy-endpoint Berezin convolution}.
\end{gathered}
\end{equation*}
This route makes the data norm and the exceptional edge $n=|m|$ visible rather than concealing them in a formal block sum.

Several related questions remain outside the scope of the present paper.  Frequency-weighted DtN or NtD data lead to logarithmic, dilogarithmic or shifted-polylogarithmic kernels; for the dilogarithm kernel a third-order adjoint intertwiner is available, but no explicit singular system is presently known.  If the perturbation is known near the boundary, the radial operator becomes compact and a commuting Heun equation appears.  These are separate spectral problems and are not needed for the present full-data result.

\section{Checks required before submission}
\label{sec:todos}

The red items in the source are deliberately retained at this stage.  The principal checks are:
\begin{enumerate}
\item derive \eqref{eq:fourier-matrix-element} from the exact convention used for the Schr\"odinger boundary map and settle all signs and constants;
\item state precisely how zero boundary frequency is handled and how the Schr\"odinger formulation relates to conductivity EIT;
\item verify that the chosen non-redundant norm is the intended equal-information norm, and distinguish it from physical noise weighting;
\item check \eqref{eq:kernel-commutator} independently by symbolic calculation;
\item verify the spherical-transform normalization from a primary Berezin or Mehler--Fock source;
\item complete the numerical convergence study or reduce the numerical claims accordingly;
\item read every cited paper and audit all bibliography metadata;
\item finalize authorship, CRediT roles, data/code archive and the generative-AI disclosure.
\end{enumerate}

\ack{This draft was developed with substantial assistance from ChatGPT (OpenAI), using GPT-5.6 Sol Pro through a ChatGPT Pro subscription, accessed in September 2026.  The assistance included proposing candidate derivations, symbolic and numerical checks, literature-search support and drafting.  The human author is responsible for checking every argument, computation, citation and originality claim.  \todo{Revise this disclosure after the final human verification record and author contributions are complete.}}

\roles{William R B Lionheart: conceptualization, methodology, investigation, validation, writing, and supervision.  \todo{Replace by the final CRediT statement if further authors join the paper.}}

\data{The MATLAB code used for the exploratory finite-section figures is included with the present draft package.  A versioned public repository and a precise computational environment will be supplied with the submitted manuscript. \todo{Add repository DOI or URL and software versions.}}

\bibliographystyle{iopart-num}
\bibliography{references}

\appendix
\section{A direct polar-coordinate check of the kernel identity}
\label{app:polar-check}

Set
\[
 \psi=\theta-\phi,\qquad
 D=1-2r\rho\cos\psi+r^2\rho^2,
 \qquad \kappa=(\pi D)^{-1}.
\]
The operator acting in the $z$ variable is
\[
 \mathcal P_z=
 \frac1{4r}\partial_r\bigl(r(1-r^2)^2\partial_r\bigr)
 +\frac14(r^2+r^{-2})\partial_\psi^2+r^2,
\]
and $\mathcal P_w$ is obtained by replacing $r$ by $\rho$.  Expansion over the common denominator $\pi D^3$ gives a numerator divisible by $r^2-\rho^2$ from the radial terms; the angular and zeroth-order terms cancel the remainder, leaving
\[
 (\mathcal P_z-\mathcal P_w)\kappa=0.
\]
For reproducibility, the intended independent Mathematica check is:
\begin{lstlisting}[language=Mathematica,basicstyle=\ttfamily\footnotesize,frame=single,breaklines=true]
ClearAll[r, rho, psi, ker, Pz, Pw];
den = 1 - 2 r rho Cos[psi] + r^2 rho^2;
ker = 1/(Pi den);
Pz[u_] := 1/(4 r) D[r (1-r^2)^2 D[u,r],r]
          + 1/4 (r^2 + 1/r^2) D[u,{psi,2}] + r^2 u;
Pw[u_] := 1/(4 rho) D[rho (1-rho^2)^2 D[u,rho],rho]
          + 1/4 (rho^2 + 1/rho^2) D[u,{psi,2}] + rho^2 u;
FullSimplify[Pz[ker]-Pw[ker],
 Assumptions -> 0<r<1 && 0<rho<1 && Element[psi,Reals]]
\end{lstlisting}

\section{Conventions and alternative normalizations}
\label{app:conventions}

Changing from normalized boundary modes $(2\pi)^{-1/2}e^{ik\theta}$ to unnormalized modes rescales the data and hence $K$ by an overall constant.  Reversing the sign convention in the Fourier coefficient changes $m$ to $-m$ but leaves all kernels and spectra unchanged.  Replacing $(-\Delta+q)u=0$ by $\Delta u+qu=0$ reverses the sign of the linearized boundary map but not its normal operator.  These harmless choices should be separated from two non-harmless changes:
\begin{enumerate}
\item excluding or quotienting the constant boundary mode, which can alter low-rank edge terms;
\item weighting measurements by powers of $|k|$ and $|l|$, which changes the radial coefficient sequence and therefore the degree of ill-posedness.
\end{enumerate}

\end{document}
